\chapter{Desarrollo del trabajo}

Este capítulo no documenta el sistema, sino el proceso por el que se
construyó. La observación que lo recorre es que trabajar de forma
habitual con \textit{coding agents} --- en este caso, Claude Code ---
desplaza el centro de gravedad del ciclo clásico de desarrollo de
software: la implementación deja de ser la actividad dominante y, en su
lugar, lo son la especificación, la revisión y la verificación. La
sección que sigue desarrolla esta idea; el resto del capítulo describe el
bucle de trabajo resultante y cómo se manifestó en esta segunda fase del
proyecto.

\section{El problema con el ciclo clásico}

Los modelos de proceso clásicos --- del ciclo en cascada a sus
variantes ágiles --- tienden a asumir que el coste dominante del
desarrollo es escribir y depurar el código. Sobre esa premisa se
construye buena parte de las prácticas habituales: separar análisis de
implementación para no malgastar trabajo, escribir documentación
intermedia para guiar al programador, planificar \textit{sprints} en
función de la velocidad de implementación.

Cuando los modelos de lenguaje actuales actúan como ejecutores
principales, esa premisa queda cuestionada. En tareas generales de
programación, escribir una función concreta o un módulo completo a
partir de una especificación clara deja de ser la parte costosa del
trabajo. La evidencia empírica de este desplazamiento --- los
experimentos de Peng et al.\ con GitHub Copilot \cite{peng2023impact} y
los ensayos aleatorizados a gran escala de Cui et al.\
\cite{cui2025effects} --- se discutió en el capítulo de estado del arte
y no se repite aquí.

El peso se desplaza, en consecuencia, hacia tres actividades que antes
estaban distribuidas o asumidas: enunciar bien la especificación,
revisar lo que el agente ha producido y verificar que sigue
funcionando lo que ya funcionaba. El ciclo deja de funcionar de forma
estrictamente secuencial y se aproxima a un bucle corto con cuatro
pasos --- especificar, ejecutar, revisar, ajustar --- que se repite
muchas veces al día. En la práctica, especificar deja de ser
\emph{escribir un párrafo de requisitos} y pasa a ser \emph{redactar un
plan de implementación}: lo bastante rico como para que el agente
pueda ejecutarlo sin tener que interpretar el contrato cada pocos
pasos.

\section{El nuevo bucle de desarrollo}

El bucle adoptado durante esta fase fue este:

\begin{enumerate}
\item \textbf{Especificar.} Cada tarea no trivial empezó con una
\textit{spec} en Markdown que, más que un documento de requisitos, era
un plan de implementación por fases. Cada fase declaraba qué
pruebas debían pasar al final, qué archivos se tocaban y qué
invariantes debían preservarse. El objetivo era doble: que el agente
pudiera ejecutar fase por fase de forma incremental y segura, y que el
revisor humano dispusiera de un contrato \textit{ejecutable mentalmente}
sobre el que comparar el resultado.
\item \textbf{Ejecutar.} El agente se encarga de la implementación, las
pruebas asociadas y la mayor parte de la depuración local. En tareas
independientes, varios subagentes trabajan en paralelo sobre
\textit{worktrees} de Git separados, de modo que cada uno ve su propia
rama y su propio \textit{working tree} y no hay riesgo de que se pisen
entre sí al modificar los mismos ficheros. En tareas con estado
compartido se ejecuta de forma secuencial.
\item \textbf{Revisar.} Antes de fusionar, el código pasa por una
revisión combinada: una lectura humana enfocada en alcance y diseño, y
una segunda pasada por subagentes especializados de revisión
(\texttt{code-reviewer}, \texttt{code-simplifier} y
\texttt{silent-failure-hunter}) configurados para reportar solo hallazgos
con confianza alta.
\item \textbf{Ajustar.} Los hallazgos confirmados se convierten en una
nueva iteración, normalmente con su propia \textit{spec} reducida. Las
correcciones triviales se aplican en línea. Para los errores
funcionales se sigue el ciclo de TDD: primero se replica el fallo con
una prueba que efectivamente falla en el estado actual, después se
implementa el ajuste mínimo que la hace pasar y, por último, se vuelve
a ejecutar la suite completa para comprobar que el resto sigue verde.
\end{enumerate}

Este bucle, repetido decenas de veces, sustituye al cronograma rígido
del ciclo clásico. La planificación de alto nivel sigue existiendo
--- en este proyecto, en forma de archivos \texttt{spec.md} versionados
en el propio repositorio, no de \textit{tickets} externos --- pero
deja de gobernar el día a día.

\section{Reparto de responsabilidades}

La diferencia más clara con el ciclo clásico es el reparto del trabajo
entre humano y agente.

\paragraph{Lo que sigue siendo humano.} Las decisiones de alcance, las
elecciones arquitectónicas con consecuencias en otros módulos, la
aceptación de revisiones que el agente no es capaz de juzgar y la
interpretación de lo que el tutor pide cuando la petición es ambigua. El
contrato con la primera fase del proyecto, por ejemplo, se discutió y
estabilizó por completo antes de delegar nada. El diseño visual del
\textit{dashboard} también se mantuvo en el lado humano: la paleta, el
modo oscuro, la rejilla de tarjetas, la composición del panel de
agentes y la disposición del \textit{replay} se prototiparon en Figma
antes de empezar a programarlos. El agente intervino después,
traduciendo el diseño aprobado a componentes React y a clases de
Tailwind, en un ciclo iterativo en el que primero se cerraban unos
pocos componentes base y luego se construían pantallas reutilizándolos.

\paragraph{Lo que se delega.} La implementación de cada función o módulo
una vez fijado su contrato, la escritura de pruebas a partir de una
descripción del comportamiento esperado, los refactores mecánicos y la
depuración local que no requiere contexto del producto. La mayor parte
del código de esta segunda fase pasó por este reparto.

Para que esa delegación produzca código de calidad sin necesidad de
una corrección posterior intensiva, las \textit{specs} se redactaron en
clave de TDD: cada fase del plan declara explícitamente las pruebas
que el agente debe hacer pasar, indicando incluso las rutas concretas
de los archivos de prueba que se esperan. El agente implementa primero
esas pruebas (o las completa si están parcialmente escritas), comprueba
que fallan en el estado actual y solo entonces escribe el código que
las hace pasar. El diseño del plan es, por tanto, en gran parte el
diseño de su propia suite de validación.

Una intuición útil que emergió durante el trabajo es la siguiente: el
agente es excelente cuando se le da un contrato preciso, y poco fiable
cuando se le pide que infiera el contrato a partir del contexto. La
mayoría de los errores de juicio aparecieron cuando se intentó saltar
ese paso por economía de tiempo.

El caso más claro fue la integración inicial de las herramientas en el
Orchestrator. La primera versión daba por hecho que toda llamada que
llegara del modelo traería el campo \texttt{query} bien formado, y por
eso accedía a él directamente sobre el diccionario de parámetros. La
suposición era cómoda pero ingenua: con que el modelo serializara mal
una sola llamada o se la dejara a medias, el acceso reventaba y un
fallo \emph{de él} pasaba a ser un error \emph{del backend}. Al
revisarlo, se sustituyó el acceso directo por una lectura con valor por
defecto, se añadieron pruebas para llamadas malformadas y se dejó la
herramienta fallando con un mensaje controlado en lugar de propagar la
excepción. La conclusión no fue que el agente no supiera escribir esa
pieza --- la escribió en pocos segundos cuando se le pidió ---, sino
que el contrato que se le dio no estaba lo suficientemente cerrado como
para que pudiera anticipar el caso degenerado.

\section{Planificación y coordinación con la primera fase}

Esta fase se desarrolló en paralelo con la primera. Para evitar
acoplamiento prematuro entre ambas se siguió una estrategia explícita:
fijar pronto el contrato de integración y dejar que cada fase evolucionara
por su cuenta sobre ese contrato. El
contrato resultante --- los tres métodos \texttt{decide},
\texttt{update} y \texttt{get\_state} --- es lo suficientemente mínimo
como para que ninguna de las dos fases haya necesitado coordinar
\textit{releases} con la otra durante los meses siguientes.

La coordinación entre fases no se apoyó en un gestor de \textit{tickets}
externo, sino en el propio repositorio compartido en GitHub. Cada
cambio relevante --- de cualquiera de las dos fases o de la capa común
--- venía acompañado de un archivo \texttt{spec.md} versionado en la
misma rama, que funcionaba a la vez como plan para el agente y como
documento de discusión con el otro autor. Los cambios viajaban por las
ramas del repositorio: cuando uno de los dos hacía \texttt{pull}, veía
de forma automática tanto el código nuevo como las \textit{specs} que
lo habían motivado. La comunicación complementaria del día a día se
mantuvo en mensajería directa, reservada a desbloqueos puntuales. Los
cambios que afectaban al esquema compartido del Knowledge Backbone se
discutían en la \textit{spec} correspondiente antes de implementarse en
cualquiera de las dos fases.

\section{Hitos de iteración del sistema}

El sistema final no se planificó como un todo, sino que emergió por
acumulación de iteraciones. Las más importantes, en orden
aproximado, fueron las siguientes:

\paragraph{M1. Pipeline básico sobre CLI.} La primera versión funcional
se operaba enteramente desde una terminal: el usuario describía el
problema y el Orchestrator invocaba en secuencia al Architect, al
Tracker, al Analyst y al Reporter, sin pausas intermedias ni
recomendaciones al usuario; era un \textit{pipeline} de un solo tirón
cuya salida se imprimía en la consola. No había visualización del
entorno ni persistencia entre sesiones. El objetivo de esta versión era
validar que los cinco roles podían coordinarse sin solaparse, no
producir una interfaz de uso.

\paragraph{M2. Interfaz web con \textit{replay}.} Una vez estabilizado
el pipeline, se añadió un \textit{frontend} en React conectado al
backend mediante WebSocket. El cambio principal no fue cosmético sino
de paradigma: cada paso del pipeline empezó a publicarse como un evento
observable desde la interfaz, con \textit{replay} paso a paso de la
simulación y paneles de estado por agente. La CLI quedó relegada a un
canal de pruebas y depuración.

\paragraph{M3. Comparativa multi-modelo.} El \texttt{model\_loader} se
extendió para permitir instanciar varios modelos de la primera fase en
el mismo entorno. El resto del pipeline se adaptó para tratar las
trayectorias por agente en lugar de por simulación. Esta iteración
ocurrió antes de la integración con el Knowledge Backbone, y fue lo que
fijó el resto del pipeline en términos de \emph{múltiples agentes en
una misma ejecución} en lugar de \emph{una simulación, un modelo}.

\paragraph{M4. Orchestrator más conversacional.} Antes de abrir el
sistema al Knowledge Backbone, el Orchestrator pasó por una iteración
intermedia orientada al uso: dejó de ser un despachador silencioso de
\textit{tool calls} y empezó a recomendar al usuario qué hacer a
continuación, qué modelos parecían encajar con el problema descrito y
qué herramientas tenía a su disposición. Esta iteración no introdujo
nuevos agentes, pero hizo que la interfaz dejara de sentirse como un
botón de \emph{lanzar} y se pareciera más a una conversación con
sugerencias.

\paragraph{M5. Lectura del Knowledge Backbone.} El Analyst dejó de
trabajar exclusivamente sobre los datos del Tracker y empezó a consultar
postulados previos en el grafo, a través de \texttt{retrieve\_context}.
Poco después se extendió el patrón al Architect y al Reporter, pero en
su variante inyectada: el Orchestrator pre-recupera el contexto y lo
entrega como parámetro.

\paragraph{M6. Persistencia de observaciones (\textit{sim-memory}).} El
Tracker pasó de ser un observador puro a alimentar el Knowledge
Backbone, mediante \texttt{TrackerMemoryWriter}. Esta iteración introdujo
la tabla \texttt{simulation\_observations} en PostgreSQL y el
\textit{namespace} \texttt{simulation} en Qdrant, y convirtió la segunda
fase en productora del backbone además de consumidora.

\paragraph{M7. BM25 nativo en Qdrant.} La búsqueda \textit{sparse} se
migró del cliente Python a la implementación nativa de Qdrant,
delegando la tokenización y el cálculo de pesos en el servidor.

\paragraph{M8. Generación de informes por secciones.} El Reporter sufrió
una limitación práctica: los informes completos podían superar el
presupuesto de salida del modelo o producir LaTeX difícil de reparar. La
solución fue desplazar parte de la estructura al backend: el sistema pide
al modelo secciones acotadas, ensambla el documento final y conserva una
ruta de reserva determinista si alguna sección falla. Este cambio hizo
menos elegante la implementación, pero más fiable la entrega del PDF.

\section{La verificación como nuevo cuello de botella}

Si la implementación deja de ser lo costoso, el problema se desplaza:
¿qué garantiza que el código generado es correcto?

La respuesta adoptada en este proyecto tiene tres capas. La primera es
\textbf{las pruebas como contrato}: una funcionalidad sin prueba no se
considera terminada. La suite combina pruebas unitarias y de integración
en \textit{pytest} para el backend, y pruebas \textit{end-to-end} en
Playwright para el cliente web. La segunda capa es la \textbf{revisión
combinada del cambio antes de fusionar}: la propia lectura humana del
\textit{diff} --- centrada en el alcance, en lo que el cambio no
debería romper y en si lo entregado coincide con la \textit{spec} --- y
una pasada complementaria por subagentes especializados
(\texttt{code-reviewer}, \texttt{code-simplifier} y
\texttt{silent-failure-hunter}) configurados para reportar solo
problemas de alta confianza. Los subagentes no reemplazan la revisión
humana; la enfocan, encontrando antes los problemas mecánicos y
dejando a la lectura humana las cuestiones de diseño y alcance. La
tercera capa es la \textbf{verificación manual sobre la aplicación
viva}: para los cambios que afectan a la interfaz, no basta con que las
pruebas pasen; el cambio debe poderse usar en un navegador y mostrarse
de forma correcta antes de aceptarlo.

Esta capa evitó varias regresiones pequeñas pero costosas. En la
generación de informes, por ejemplo, las pruebas unitarias podían validar
que existía una ruta de compilación, pero la ejecución completa reveló
fallos distintos: respuestas truncadas del Reporter, secciones duplicadas
o mensajes de descarga emitidos dos veces en el chat. Esos errores no
eran fallos de sintaxis aislados; afectaban al recorrido real del usuario.
Por eso la verificación terminó combinando pruebas de bajo nivel con una
sesión completa desde la interfaz.

\section{Limitaciones observadas}

El bucle descrito no resuelve todo. Durante el desarrollo aparecieron tres
limitaciones recurrentes.

La primera es que el agente puede generar resultados convincentes sin
ser correctos, especialmente cuando el contrato es vago. Esto se mitiga
con \textit{specs} precisas, pero no se elimina.

La segunda es de ritmo. Cuando el coste de implementar baja, lo natural
es producir más cambios; pero el coste de \emph{revisar} no baja al
mismo ritmo. Si uno se deja llevar por la velocidad de generación,
acaba acumulando cambios mal mirados que después cuestan más de
desenredar de lo que se ahorró al hacerlos. Aquí se aplicó una regla
sencilla para no caer en eso: ningún cambio se considera terminado
hasta que ha pasado por la revisión combinada del bucle y, cuando toca
interfaz, por una comprobación a mano sobre la aplicación viva. Esa
tensión entre velocidad de producción y carga de revisión aparece
también en los estudios cualitativos sobre el uso real de asistentes
de IA en entornos profesionales \cite{butler2024dear}.

La tercera es metodológica, y afecta al propio acto de escribir esta
memoria. Documentar un sistema acabado consiste en describir un
conjunto cerrado de decisiones; documentar un proceso como este,
en cambio, exige asumir que no hubo grandes decisiones, sino muchas
pequeñas iteraciones cuyo valor solo se ve cuando se miran encadenadas.
Por eso este capítulo no presenta el desarrollo como un plan ejecutado
de principio a fin, sino como un bucle corto que se repitió una y otra
vez. Disfrazarlo de proceso lineal habría sido más cómodo de leer,
pero menos fiel a lo que de verdad pasó.

\section{Retos prácticos: la memoria de los agentes en tokens}

Una limitación que apareció pronto y que condicionó decisiones de
diseño concretas es la memoria de los propios agentes. Los modelos de
lenguaje operan sobre una ventana de contexto finita medida en
\textit{tokens}, y en una sesión real del laboratorio esa ventana se
llena con sorprendente facilidad: el historial conversacional con el
usuario, las trazas completas de la simulación, los JSON del Tracker
con eventos paso a paso y las memorias recuperadas del Knowledge
Backbone. Si todo eso se vuelca a cada agente en cada turno, el
Orchestrator empieza a perder precisión, a ignorar parte de las
instrucciones del \textit{system prompt} y a alucinar detalles que ya
no caben de forma íntegra en su contexto.

La respuesta a este problema se dio en dos frentes. Por un lado, se
introdujo un mecanismo de \textbf{autocompactación} del historial: a
partir de cierto umbral de uso de la ventana, el Orchestrator
reemplaza los mensajes más antiguos por un resumen estructurado que
preserva las decisiones tomadas, los \textit{artefactos} producidos
hasta ese punto y las preferencias del usuario, pero descarta los
detalles ya consolidados. Por otro lado, se diseñaron herramientas
explícitas de \textbf{detalle y resumen} sobre los datos voluminosos: en
lugar de inyectar el JSON completo del Tracker en cada turno, el
Orchestrator entrega de forma rutinaria solo un resumen estructurado
y deja a disposición \textit{tool calls} específicas
(\texttt{get\_tracker\_detail}, \texttt{get\_analyst\_detail},
\texttt{get\_simulation\_step\_window}, entre otras) para que el agente
recupere bajo demanda únicamente la porción de detalle que necesita en
un momento dado.

Este patrón --- resumen por defecto, detalle bajo demanda --- desplaza
parte de la complejidad al diseño de las herramientas, pero permite
trabajar con sesiones largas sin saturar el contexto del modelo.
También subraya una idea más general que aparece a lo largo del
proyecto: gran parte del trabajo con agentes consiste en decidir qué
ven y qué no ven, no en escribir el código que los invoca.
